\documentclass[aspectratio=169]{beamer}

% ─────────────────────────────────────────────
%  Theme & colours
% ─────────────────────────────────────────────
\usetheme{Madrid}
\usecolortheme{default}
\definecolor{WUBlue}{RGB}{0,62,120}
\definecolor{WUGold}{RGB}{192,145,0}
\definecolor{LightBlue}{RGB}{55,138,221}
\definecolor{DarkBlue}{RGB}{21,47,95}
\definecolor{AccentOrange}{RGB}{224,92,42}
\definecolor{AccentGreen}{RGB}{40,150,80}

\setbeamercolor{palette primary}{bg=WUBlue,fg=white}
\setbeamercolor{palette secondary}{bg=WUBlue!80,fg=white}
\setbeamercolor{palette tertiary}{bg=WUBlue!60,fg=white}
\setbeamercolor{palette quaternary}{bg=WUBlue,fg=white}
\setbeamercolor{structure}{fg=WUBlue}
\setbeamercolor{section in toc}{fg=WUBlue}
\setbeamercolor{alerted text}{fg=AccentOrange}
\setbeamercolor{block title}{bg=WUBlue,fg=white}
\setbeamercolor{block body}{bg=WUBlue!8}

\setbeamertemplate{navigation symbols}{}
\setbeamertemplate{footline}[frame number]

% ─────────────────────────────────────────────
%  Packages
% ─────────────────────────────────────────────
\usepackage{booktabs}
\usepackage{array}
\usepackage{graphicx}
\usepackage{amsmath}
\usepackage{xcolor}
\usepackage{colortbl}
\usepackage{tikz}
\usepackage{multirow}
\usepackage{tabularx}
\usepackage{adjustbox}
\usepackage{caption}

% ─────────────────────────────────────────────
%  Title page
% ─────────────────────────────────────────────
\title[GEP Index — Final Results]%
      {\textbf{Geoeconomic Pressure Index}\\[4pt]
       \large Final Results: GTM-6 Construction, MIN2 Index \& Return Predictability}
\author[Caterina Piacentini]{Caterina Piacentini}
\institute[WU Vienna]{Vienna University of Economics and Business\\
           \smallskip\small Thesis Advisor: Dangl \& Salbrechter methodology}
\date{April 2026}

% ─────────────────────────────────────────────
\begin{document}
% ─────────────────────────────────────────────

\begin{frame}
  \titlepage
\end{frame}

% ─── Outline ─────────────────────────────────
\begin{frame}{Outline}
  \tableofcontents
\end{frame}

% ════════════════════════════════════════════════════════
\section{GTM-6 — Final Topic Model}
% ════════════════════════════════════════════════════════

\begin{frame}{GTM-6: Six Geoeconomic Sub-topics (New Final)}

The \textbf{Guided Topic Model (GTM-6)} expands six seed clusters
in a 64-dimensional Word2Vec embedding space trained on
30 years of Reuters newswire (1996--2025).

\medskip
\begin{columns}[T]
\column{0.52\textwidth}
\begin{block}{Sub-topics (6) — new final}
\begin{enumerate}\small
  \item Sanctions
  \item \textbf{Trade Coercion} \textcolor{AccentOrange}{(new)}
  \item Export Controls
  \item Financial Coercion
  \item Retaliation
  \item Embargo
\end{enumerate}
\end{block}

\column{0.44\textwidth}
\begin{block}{Key parameters}
\begin{itemize}\small
  \item Cluster size: 100 words each
  \item Gravity: 1.5
  \item $\alpha_\text{max}$: 2.0 rad
  \item FAISS $k$-NN: 5\,000 candidates
\end{itemize}
\end{block}
\end{columns}

\bigskip
\begin{alertblock}{Consolidation from 8 to 6 topics}
  Trade War, Tariffs, and Protectionism cluster tightly in embedding space
  $\to$ merged into \textbf{Trade Coercion}.
  Eliminates topic fragmentation; improves semantic coherence.
\end{alertblock}

\end{frame}

% ─── Seed words table ────────────────────────
\begin{frame}{GTM-6 — Seed Words per Sub-topic}

\begin{center}
\footnotesize
\renewcommand{\arraystretch}{1.25}
\resizebox{\textwidth}{!}{%
\begin{tabular}{@{}llll@{}}
\toprule
\textbf{Sub-topic} & \textbf{Positive seeds} & \textbf{Negative seeds} & \textbf{Design note} \\
\midrule
\rowcolor{WUBlue!10}
Sanctions
  & \texttt{economic\_sanctions}, \texttt{targeted\_sanctions}
  & \texttt{sanctions\_relief}, \texttt{sanctions\_waiver}
  & Coercive intent \\
Trade Coercion
  & \texttt{trade\_war}, \texttt{retaliatory\_tariffs}
  & \texttt{trade\_deal}, \texttt{trade\_pact}
  & Merges Trade War + Tariffs + Protectionism \\
\rowcolor{WUBlue!10}
Export Controls
  & \texttt{export\_ban}, \texttt{entity\_list}
  & \texttt{export\_license}, \texttt{export\_licenses}
  & Captures US BIS Entity List / chip controls \\
Financial Coercion
  & \texttt{asset\_freeze}, \texttt{secondary\_sanctions}
  & \texttt{debt\_relief}
  & Sovereign-level financial weapons \\
\rowcolor{WUBlue!10}
Retaliation
  & \texttt{retaliation}, \texttt{countermeasures}
  & \texttt{concessions}, \texttt{goodwill\_gesture}
  & Escalation \& counter-escalation \\
Embargo
  & \texttt{trade\_embargo}, \texttt{oil\_embargo}
  & \texttt{lift\_embargo}, \texttt{lifting\_sanctions}
  & Added \texttt{lifting\_sanctions} as negative \\
\bottomrule
\end{tabular}}
\end{center}

\vspace{4pt}\footnotesize
Positive seeds \emph{anchor} the subspace; negative seeds act as
\emph{repellors}, pushing expansion away from resolution/relief language.

\end{frame}

% ─── Dictionary aggregation ──────────────────
\begin{frame}{Dictionary Construction — Aggregation Mechanism}

Each sub-topic produces 100 words ranked by \emph{projection norm} onto
the topic subspace. The global dictionary is built in two steps
(Dangl \& Salbrechter 2023, Eq.~10):

\medskip
\textbf{Step 1 — Local normalisation (per topic $q$):}
\[
  \tilde{w}_{q,j} = \frac{w_{q,j}}{\max_j\, w_{q,j}} \;\in\; [0,1]
\]

\textbf{Step 2 — Global average (across $Q=6$ topics):}
\[
  \bar{w}_j = \frac{1}{Q}\sum_{q=1}^{Q} \tilde{w}_{q,j}
\]

\medskip
Words appearing in \emph{multiple} sub-topics receive proportionally
higher global weights, reflecting cross-dimensional geoeconomic relevance.

\medskip
\begin{block}{Final dictionary (new final, $Q=6$)}
  6 sub-topics $\times$ 100 words each.
  Top entries span all sub-dimensions:
  \texttt{trade\_war}, \texttt{retaliatory\_tariffs}, \texttt{sanctions},
  \texttt{entity\_list}, \texttt{retaliation}, \texttt{embargo}\ldots
\end{block}

\end{frame}

% ─── Word clouds ────────────────────────────
\begin{frame}{Word Clouds — All 6 Sub-topics (New Final)}
  \begin{center}
    \includegraphics[width=0.88\textwidth]%
      {GTM_different_versions/GTM_new_final_results/Combined_GEP_Grid_6.png}
  \end{center}
\end{frame}

% ════════════════════════════════════════════════════════
\section{MIN2 Index — Construction}
% ════════════════════════════════════════════════════════

\begin{frame}{GEP Index — From Scoring to Threshold Classification}

\textbf{Key innovation:} the final index uses a \emph{threshold-based}
classification instead of continuous article scoring.

\medskip
\begin{columns}[T]
\column{0.48\textwidth}
\begin{block}{Old approach (score-based)}
\small
Per-article score:
\[
  s_i = \frac{\sum_w \min(\text{cnt}_w,\,4)\cdot T(w)}{N_i}
\]
GEP$_m$ = article-weighted mean of $s_i$.

\smallskip
Problems:
\begin{itemize}
  \item Scale $\sim 10^{-4}$ (uninterpretable)
  \item \textbf{Non-stationary in levels}\\
        $\Rightarrow$ must use $\Delta$GEP
  \item Sensitive to single high-freq word
\end{itemize}
\end{block}

\column{0.48\textwidth}
\begin{block}{New approach — MIN-$K$ threshold}
\small
Article is a GEP article if it matches $\geq K$ dictionary keywords.
\[
  \text{GEP}_m = \frac{n_{\text{GEP articles},\,m}}{n_{\text{articles},\,m}}
\]

Properties:
\begin{itemize}
  \item Proportion in $[0,1]$ — share of news devoted to geoeconomic pressure
  \item \textcolor{AccentGreen}{\textbf{Stationary in levels}} (ADF confirmed)
  \item Co-occurrence requirement filters noise
  \item Can use GEP$_z$ \emph{levels} as regressor
\end{itemize}
\end{block}
\end{columns}

\end{frame}

% ─── MIN-K variants ──────────────────────────
\begin{frame}{MIN-K Robustness Variants — Choosing MIN2 as Final}

\begin{center}
\footnotesize
\renewcommand{\arraystretch}{1.3}
\begin{tabular}{@{}clll@{}}
\toprule
\textbf{Variant} & \textbf{Threshold $K$} & \textbf{Interpretation} & \textbf{Status} \\
\midrule
\rowcolor{WUBlue!5}
MIN1 & $\geq 1$ keyword & Any mention of a GEP word & Robustness \\
\rowcolor{WUBlue!20}
\textbf{MIN2} & $\boldsymbol{\geq 2}$ \textbf{keywords} & \textbf{Co-occurrence required} & \textbf{FINAL} \\
\rowcolor{WUBlue!5}
MIN3 & $\geq 3$ keywords & Substantially GEP-focused article & Robustness \\
MIN4 & $\geq 4$ keywords & Most conservative & Robustness \\
\bottomrule
\end{tabular}
\end{center}

\bigskip
\begin{columns}[T]
\column{0.48\textwidth}
\begin{block}{Why MIN2?}
\begin{itemize}\small
  \item MIN1 is noisy — a single word (e.g., \texttt{tariff}) can
        appear incidentally in unrelated stories
  \item MIN2 requires \emph{co-occurrence} — the article must
        discuss geoeconomic pressure substantively
  \item MIN3/MIN4 are too conservative, losing relevant episodes
\end{itemize}
\end{block}

\column{0.48\textwidth}
\begin{block}{Key MIN2 peaks (GEP\_monthly)}
\small
\begin{tabular}{ll}
\toprule
\textbf{Event} & \textbf{GEP} \\
\midrule
9/11 (2001-09) & 0.129 \\
Iraq War (2003-02) & 0.122 \\
US--China tariffs (2018-05) & 0.107 \\
Ukraine invasion (2022-03) & 0.139 \\
Liberation Day (2025-04) & \textbf{0.197} \\
\bottomrule
\end{tabular}
\end{block}
\end{columns}

\end{frame}

% ════════════════════════════════════════════════════════
\section{Index Plots \& Validation}
% ════════════════════════════════════════════════════════

\begin{frame}{MIN2 GEP Index — Monthly (1996–2025)}
  \begin{center}
    \includegraphics[width=0.95\textwidth]%
      {INDEX/index_new_final/MIN2/GEP_Monthly_Robust_min2.png}
  \end{center}
  \vspace{-4pt}
  \footnotesize
  Share of monthly articles matching $\geq 2$ dictionary keywords.
  Annotated peaks correspond to key geoeconomic episodes.
\end{frame}

\begin{frame}{MIN2 GEP Index — Normalised (avg = 100)}
  \begin{center}
    \includegraphics[width=0.95\textwidth]%
      {INDEX/index_new_final/MIN2/GEP_Monthly_Robust_min2_norm100.png}
  \end{center}
  \vspace{-4pt}
  \footnotesize
  EPU-style normalisation: long-run average = 100.
  Liberation Day (April 2025) registers at $\sim$200 — twice the historical average.
\end{frame}

\begin{frame}{GEP vs.\ GPR (Caldara \& Iacoviello 2022)}

\textbf{External benchmark:} GPR captures \emph{geopolitical} risk;
MIN2 GEP captures \emph{geoeconomic} pressure.
Both series z-scored before comparison.

  \begin{center}
    \includegraphics[width=0.92\textwidth]%
      {INDEX/index_new_final/MIN2/gep_vs_gprd_timeseries.png}
  \end{center}

\vspace{-4pt}
\footnotesize
Low but positive Pearson $r$ —
GEP and GPR are \emph{related but distinct}: they diverge most
during purely economic coercion episodes (trade war 2018--2019)
and converge during joint military--economic escalation (Ukraine 2022).

\end{frame}

\begin{frame}{Cross-Correlation: MIN2 GEP vs.\ GPR}
  \begin{center}
    \includegraphics[width=0.88\textwidth]%
      {INDEX/index_new_final/MIN2/gep_vs_gprd_crosscorr.png}
  \end{center}
  \medskip\small
  Peak correlation at positive lags confirms that geoeconomic pressure
  (trade tensions, sanctions escalation) tends to \emph{precede}
  escalation in geopolitical risk readings.
\end{frame}

% ════════════════════════════════════════════════════════
\section{Return Predictability Regressions}
% ════════════════════════════════════════════════════════

\begin{frame}{Regression Setup}

\textbf{Dependent variable:} S\&P 500 monthly (and quarterly) log return.

\textbf{Key regressor:} GEP$_z$ — z-scored MIN2 index \textbf{in levels}.
Using levels is valid because the bounded [0,1] ratio is stationary (ADF confirmed).

\medskip
\begin{center}
\footnotesize
\renewcommand{\arraystretch}{1.2}
\begin{tabular}{@{}cll@{}}
\toprule
\textbf{Spec} & \textbf{Formula} & \textbf{Purpose} \\
\midrule
m1 & $r_t = \alpha + \beta\,\text{GEP}_z + \varepsilon$ & Univariate baseline \\
m2 & $+ \text{GPR}_z$ & Separates GEP from geopolitical risk \\
m3 & $+ \text{MktRF} + \text{SMB} + \text{HML}$ & Controls for factor exposure \\
m4 & $+ \text{GEP}_{z,t-1} + \text{GPR}_{z,t-1}$ & Lag structure \\
m\_rob & $r_t = \alpha + \beta\,\Delta\text{GEP}_z + \gamma\,\text{GPR}_z + \varepsilon$ & Robustness: differences \\
\bottomrule
\end{tabular}
\end{center}

\medskip
\begin{itemize}
  \item \textbf{Standard errors:} Newey-West HAC (bandwidth $= \lfloor 4(T/100)^{2/9}\rfloor$)
  \item \textbf{Horizons:} $h=0$ (contemporaneous), $h=1$ (next-period predictability)
  \item \textbf{Subperiods:} Full sample, Pre-GFC (1996--2007), GFC (2008--2011),
        Post-GFC (2012--2021), Russia--Ukraine (2022--2023), 2025
\end{itemize}

\end{frame}

\begin{frame}{Main Results — Monthly, $h=0$ (Contemporaneous)}

\begin{center}
\footnotesize
\renewcommand{\arraystretch}{1.25}
\resizebox{\textwidth}{!}{%
\begin{tabular}{@{}lcccccc@{}}
\toprule
\textbf{Period (N)} & \multicolumn{2}{c}{\textbf{m1}} & \multicolumn{2}{c}{\textbf{m2}} & \textbf{m3 (GEP$_z$)} & \textbf{m\_rob ($\Delta$GEP$_z$)} \\
                    & Coef & $p$ & Coef & $p$ & Coef / $p$ & Coef / $p$ \\
\midrule
\rowcolor{WUBlue!8}
Full sample (N$\approx$355) & $-$0.005 & 0.16 & $-$0.004 & 0.25 & $+$0.001* / 0.005 & n.s. \\
Pre-GFC (N=142) & \textcolor{AccentOrange}{$-$0.009*} & 0.025 & $-$0.010. & 0.070 & n.s. & \textcolor{AccentOrange}{$-$0.017* / 0.021} \\
\rowcolor{WUBlue!8}
GFC 2008--11 (N=48) & n.s. & — & n.s. & — & n.s. & n.s. \\
Post-GFC (N=119) & n.s. & 0.76 & n.s. & 0.90 & \textcolor{AccentGreen}{$+$0.001** / 0.004} & n.s. \\
\rowcolor{WUBlue!8}
Russia--Ukraine (N=23) & n.s. & 0.23 & \textcolor{AccentOrange}{$-$0.047* / 0.021} & & n.s. & n.s. \\
\bottomrule
\end{tabular}}
\end{center}

\medskip\small
\begin{itemize}
  \item \textcolor{AccentOrange}{\textbf{Orange}}: negative — geo pressure \emph{hurts} returns (fear regime)
  \item \textcolor{AccentGreen}{\textbf{Green}}: positive — geo pressure commands \emph{risk premium} (Post-GFC)
  \item \textbf{m3 full sample}: significant only after stripping factor exposure
\end{itemize}

\end{frame}

\begin{frame}{Regime Flip: Pre-GFC vs.\ Post-GFC}

\begin{columns}[T]
\column{0.48\textwidth}
\begin{block}{\textcolor{AccentOrange}{Pre-GFC (1996--2007): Fear regime}}
\begin{itemize}\small
  \item GEP$_z$ = $-0.009^*$ in m1 (monthly)
  \item GEP$_z$ = $-0.028^*$ in m1 (\emph{quarterly} — stronger)
  \item m\_rob ($\Delta$GEP$_z$): $-0.017^*$ — levels and differences agree
  \item Effect disappears once MktRF added (m3):
        \emph{worked through broad market repricing}
  \item \textbf{Mechanism:} geopolitical news caused immediate, market-wide fear selloffs
\end{itemize}
\end{block}

\column{0.48\textwidth}
\begin{block}{\textcolor{AccentGreen}{Post-GFC (2012--2021): Risk premium regime}}
\begin{itemize}\small
  \item GEP$_z$ = $+0.0008^{**}$ in m3 (controls required)
  \item GPR$_z$ = $-0.013^*$ in m2 simultaneously
  \item GEP and GPR load with \emph{opposite signs} — they capture distinct channels
  \item \textbf{Mechanism:} in the low-vol QE era, geoeconomic uncertainty was a \emph{priced} factor
        $\to$ high-GEP months earned higher risk-adjusted returns
\end{itemize}
\end{block}
\end{columns}

\bigskip
\begin{alertblock}{Key takeaway}
  The sign of the GEP--return relationship is not stable over time.
  Fear and risk-premium regimes alternate. Structural break analysis
  (Bai-Perron) is consistent with this pattern.
\end{alertblock}

\end{frame}

\begin{frame}{Predictability at $h=1$: No Leading Effect}

\begin{center}
\footnotesize
\renewcommand{\arraystretch}{1.25}
\begin{tabular}{@{}lcccc@{}}
\toprule
\textbf{Period} & \textbf{m1} ($p$) & \textbf{m2} ($p$) & \textbf{m3} ($p$) & \textbf{m\_rob} ($p$) \\
\midrule
Full sample    & 0.70 & 0.73 & 0.54 & 0.71 \\
Pre-GFC        & 0.91 & 0.42 & 0.51 & 0.34 \\
GFC 2008--11   & 0.66 & 0.57 & 0.45 & 0.82 \\
Post-GFC       & 0.73 & 0.93 & 0.94 & 0.84 \\
Russia--Ukraine & 0.97 & 0.40 & 0.44 & 0.29 \\
\bottomrule
\end{tabular}
\end{center}

\medskip
\begin{block}{Interpretation}
\textbf{GEP is a coincident indicator, not a leading one.}
No specification delivers significant next-month (or next-quarter) return
predictability. The market processes geoeconomic information
\emph{contemporaneously} — the signal is already in prices by month-end.
\end{block}

\medskip\small
Exception: quarterly full-sample m4 has GEP$_{z,t-1}$ = $+0.0024^*$ ($p=0.044$) —
a weak quarterly lag that may reflect slow institutional portfolio rebalancing.

\end{frame}

\begin{frame}{Quarterly Results — Selected Highlights}

\begin{center}
\footnotesize
\renewcommand{\arraystretch}{1.25}
\resizebox{0.92\textwidth}{!}{%
\begin{tabular}{@{}lccccc@{}}
\toprule
\textbf{Period (N)} & \textbf{Spec} & \textbf{GEP$_z$ coef} & \textbf{$p$} & \textbf{R$^2$} & \textbf{Note} \\
\midrule
Full sample (118) & m3 & $+0.0020$ & 0.063 & 0.989 & Marginal positive \\
Full sample (118) & m4 & GEP lag $+0.0024^*$ & 0.044 & 0.989 & Lagged risk premium \\
\rowcolor{WUBlue!8}
Pre-GFC (46)      & m1 & $-0.0284^*$ & 0.022 & 0.085 & \textbf{Strong negative} \\
Pre-GFC (46)      & m3 & n.s. & — & 0.987 & FF3 absorbs \\
Post-GFC (40)     & m3 & $+0.0017^.$ & 0.052 & 0.993 & Marginal positive \\
\rowcolor{WUBlue!8}
Post-GFC (40)     & m4 & GPR lag $-0.049^*$ & 0.019 & 0.993 & GPR lag drives next quarter \\
\bottomrule
\end{tabular}}
\end{center}

\medskip\small
Quarterly aggregation (returns summed, GEP averaged within quarter)
amplifies the Pre-GFC negative signal: coefficient triples in magnitude ($-0.009 \to -0.028$),
confirming the effect is genuine and not a monthly artefact.

\end{frame}

% ════════════════════════════════════════════════════════
\section{Summary}
% ════════════════════════════════════════════════════════

\begin{frame}{Summary}

\begin{columns}[T]
\column{0.50\textwidth}
\begin{block}{Methodology (completed)}
\begin{itemize}\small
  \item GTM-6 final topic model — 6 coherent sub-topics
        (Trade Coercion consolidates 3 old topics)
  \item MIN2 index: GEP = share of articles with $\geq 2$ keywords
  \item Stationary by construction $\to$ levels-based regression
  \item Robustness: MIN1, MIN3, MIN4 variants;
        $\Delta$GEP$_z$ column in regressions
  \item External validation vs.\ GPR: low $r$,
        GEP precedes GPR in cross-correlation
\end{itemize}
\end{block}

\column{0.46\textwidth}
\begin{block}{Empirical findings (completed)}
\begin{itemize}\small
  \item \textbf{Pre-GFC:} GEP$_z$ negative and significant —
        geoeconomic pressure hurts returns (fear regime)
  \item \textbf{Post-GFC:} positive conditional on FF3 —
        geopolitical risk premium
  \item \textbf{Russia--Ukraine:} strong negative —
        fear regime returns
  \item \textbf{$h=1$:} no predictability anywhere —
        coincident, not leading indicator
  \item GEP and GPR have \emph{opposite} signs in Post-GFC:
        distinct asset pricing channels
\end{itemize}
\end{block}
\end{columns}

\bigskip
\begin{alertblock}{Central finding}
  The MIN2 GEP index captures a \textbf{regime-dependent} relationship
  with equity returns: a fear channel dominates in volatile periods
  (Pre-GFC, Ukraine) while a risk-premium channel dominates in
  low-volatility periods (Post-GFC).
  The index is \emph{distinct from GPR} and adds new information
  about geoeconomic coercion specifically.
\end{alertblock}

\end{frame}

% ─────────────────────────────────────────────
\end{document}
